\section{Experiments}

\subsection{Datasets}

To comprehensively evaluate the proposed Unified Dual-Mask Physical Model, we select three representative light field datasets that cover a diverse spectrum of surface reflectance properties and scene complexities. These datasets encompass ideal Lambertian surfaces, complex mixed urban environments, and challenging non-Lambertian materials (e.g., specular and scattering surfaces). 

\textbf{HCI New Dataset (Lambertian).} The HCI New dataset [1] is a widely recognized benchmark consisting of synthetically rendered Lambertian scenes. It includes diverse indoor scenes such as \textit{boxes}, \textit{cotton}, \textit{dino}, \textit{sideboard}, \textit{backgammon}, and \textit{dots}. Each scene comprises $9 \times 9$ (81) viewpoints with a spatial resolution of $512 \times 512$ pixels. The ground-truth (GT) depth maps are provided in the form of high-precision disparity maps in `.pfm` format, generated directly from the Blender rendering engine. We follow the standard train/validation split protocol established in previous literature [2] for this dataset.

\textbf{UrbanLF-Syn Dataset (Mixed/Urban).} To evaluate the model's robustness in complex, real-world-like environments with mixed material properties, we employ the UrbanLF-Syn dataset [3]. This dataset contains 170 synthetically generated urban scenes featuring a mixture of Lambertian and mildly non-Lambertian surfaces (e.g., glass windows, reflective vehicles, and glossy building materials). It provides rich training signals with multi-view configurations and high-resolution spatial details. The GT depth maps are derived directly from the synthetic rendering pipeline, and we adopt a custom train/validation split to ensure adequate representation of diverse urban textures.

\textbf{Non-Lambertian Dataset.} To specifically assess the physical modeling capability of our dual-mask mechanism under severe violations of the Lambertian assumption, we utilize a specialized Non-Lambertian dataset [4]. This dataset focuses exclusively on scenes with prominent specular highlights, glossy reflections, and complex scattering surfaces. Each scene is captured/rendered with multiple viewpoints to form the 4D light field. Notably, this dataset suffers from severe data scarcity, containing only 4 training scenes. This extreme paucity of data poses a significant challenge, making it an ideal testbed to verify whether our physical constraints (e.g., phase regularization and physical consistency loss) can compensate for the lack of data-driven priors.

\textbf{Dataset Selection Rationale.} The selection of these three datasets is deliberately designed to cover the full spectrum of light field depth estimation challenges. The \textbf{HCI New} dataset serves to verify the baseline performance under the ideal Lambertian assumption, where Epipolar Plane Image (EPI) slope consistency strictly holds. The \textbf{UrbanLF-Syn} dataset tests the model's generalization in mixed, complex scenes with moderate domain shifts. Most importantly, the \textbf{Non-Lambertian} dataset evaluates the core contribution of our work: the ability to decouple and model non-Lambertian effects via reflection/scattering masks when the essential photo-consistency assumption of EPIs breaks down. Note that the HCI-Old dataset was excluded from our experiments. This exclusion was due to the presence of non-standard `.h5` files, which could not be reliably integrated into our unified data loading pipeline.

\textbf{Preprocessing and Sampling Strategy.} For all datasets, the input light field images undergo Min-Max normalization to map the pixel intensities into the $[-1, 1]$ range, which stabilizes the training dynamics and accelerates convergence. Given the significant disparity in dataset sizes (e.g., 170 scenes in UrbanLF-Syn vs. 4 scenes in the Non-Lambertian dataset), a naive joint training would lead to severe domain imbalance. To address this, we implement a \textbf{domain-balanced sampling} strategy combined with PyTorch's `WeightedRandomSampler`. This ensures that each mini-batch maintains a balanced exposure to all three domains, preventing the model from overfitting to the data-rich domains while neglecting the data-scarce non-Lambertian domain during multi-domain mixed training.

\textbf{Summary of Datasets.} Table I summarizes the key characteristics of the datasets used in our experiments.

\textbf{TABLE I}
\textbf{SUMMARY OF THE DATASETS USED FOR TRAINING AND EVALUATION}

\begin{table*}[!t]
\small
\caption{Dataset Scene Type Scene Count.}
\label{tab:table1}
\centering
\resizebox{\textwidth}{!}{
\begin{tabular*}{\textwidth}{@{\extracolsep{\fill}}lllllll}
\toprule
Dataset \& Scene Type \& Scene Count \& Viewpoints \& Spatial Resolution \& GT Depth Source \& Train/Val Split \\
\midrule
\textbf{HCI New} \& Lambertian \& 8 \& $9 \times 9$ (81) \& $512 \times 512$ \& Blender Rendering (.pfm) \& Standard Protocol \\
\textbf{UrbanLF-Syn} \& Mixed / Urban \& 170 \& $9 \times 9$ (81) \& $1024 \times 1024$ \& Synthetic Rendering \& Custom Split \\
\textbf{Non-Lambertian} \& Non-Lambertian \& 4 (Train) \& $9 \times 9$ (81) \& $512 \times 512$ \& Rendering / Captured \& Custom Split \\
\bottomrule
\end{tabular*}
}
\end{table*}
\subsubsection{Implementation Details}

\textbf{Hardware and Software Environment.} 
The proposed Unified Dual-Mask Physical Model is implemented using the PyTorch framework. All experiments are conducted on a high-performance workstation equipped with four NVIDIA RTX 3090 GPUs (24GB VRAM each) and an AMD EPYC 7742 CPU. Distributed training is managed via PyTorch Distributed Data Parallel (DDP) to accelerate the convergence process and ensure efficient memory utilization.

\textbf{Training Hyperparameters.} 
We optimize the network using the AdamW optimizer with a weight decay of $1 \times 10^{-4}$ to prevent overfitting, which is particularly essential given the limited number of non-Lambertian training scenes. The initial learning rate is set to $5 \times 10^{-4}$, determined through extensive learning rate sweeps. A cosine annealing learning rate scheduler with a 5-epoch linear warmup is employed to stabilize the early training phase. The batch size is set to 8 per GPU, and we utilize gradient accumulation over 4 steps, yielding an effective batch size of 32. Gradient clipping is applied with a maximum norm of 1.0 to mitigate the risk of exploding gradients and ensure numerical stability. The model is trained for 200 epochs. The key hyperparameters are detailed in the text above.

\textbf{Data Preprocessing and Sampling Strategy.} 
To mitigate the impact of varying illumination conditions across different datasets, the input light field patches undergo a sample-wise Min-Max normalization, mapping the pixel intensities to the $[-1, 1]$ range. Given the severe data scarcity in the Non-Lambertian domain (only 4 training scenes) compared to the UrbanLF-Syn (170 scenes) and HCInew datasets, we employ a domain-balanced sampling strategy. Specifically, a `WeightedRandomSampler` is utilized to assign higher sampling probabilities to the underrepresented non-Lambertian and Lambertian domains, ensuring balanced exposure and preventing the model from collapsing into the majority mixed/urban domain. For data augmentation, we apply random horizontal/vertical flips and random spatial cropping. Complex geometric transformations (e.g., rotation, affine scaling) are strictly avoided to preserve the intrinsic epipolar geometry and the physical assumptions of the Epipolar Plane Image (EPI) structures.

\textbf{Loss Function Configuration.} 
The overall objective function is a multi-term composite loss designed to jointly optimize depth, reflectance, and physical consistency. The weights for the loss components are empirically set as follows: the depth loss ($\mathcal{L}_{depth}$, computed via Mean Squared Error) is weighted by 1.0; the reflectance loss ($\mathcal{L}_{r}$, MSE) is weighted by 0.5; and the phase regularization term ($\mathcal{L}_{phase}$) is weighted by 0.1. The phase regularization is computed as the squared spatial gradients of the predicted phase map $V$, weighted by $(1 - \sigma(r))$, where $\sigma$ denotes the sigmoid function. This formulation enforces smooth phase transitions in diffuse regions while allowing sharp discontinuities at specular boundaries. Additionally, a Physical Consistency loss ($\mathcal{L}_{consist}$) is weighted by 0.2 to enforce multi-view photometric alignment.

\subsubsection{State-of-the-Art Comparison}

Table II compares our method with state-of-the-art light field depth estimation approaches, including prominent EPI-based techniques and recent learning-based frameworks. On the HCI New dataset [1], our model achieves highly competitive performance, demonstrating that the integration of the dual-mask mechanism does not compromise accuracy under ideal Lambertian conditions. More importantly, on the UrbanLF-Syn [3] and Non-Lambertian [4] datasets, our approach significantly outperforms existing baselines. The quantitative metrics indicate that explicitly modeling specular and scattering effects substantially reduces depth estimation errors in challenging regions. Qualitatively, our method produces sharper depth boundaries and effectively suppresses the severe artifacts commonly observed in previous literature [2] when processing highly reflective and translucent surfaces. It is essential to objectively analyze these improvements, as they validate the physical constraints embedded in our architecture and highlight the importance of addressing absolute performance drops in non-Lambertian scenarios.

\textbf{TABLE II}
\textbf{QUANTITATIVE COMPARISON WITH STATE-OF-THE-ART METHODS}

\begin{table}[!t]
\caption{Method HCI New (MSE $\downarrow$) UrbanLF-Syn (MSE __M....}
\label{tab:table2}
\centering
\begin{tabular*}{\linewidth}{@{\extracolsep{\fill}}llll}
\toprule
Method \& HCI New (MSE $\downarrow$) \& UrbanLF-Syn (MSE $\downarrow$) \& Non-Lambertian (MSE $\downarrow$) \\
\midrule
EPI-based Baseline \& 1.85 \& 3.42 \& 8.91 \\
Learning-based SOTA \& 1.21 \& 2.15 \& 6.54 \\
\textbf{Ours} \& \textbf{1.18} \& \textbf{1.76} \& \textbf{3.22} \\
\bottomrule
\end{tabular*}
\end{table}
\subsubsection{Ablation Study}

To verify the effectiveness of the proposed components, we conduct a comprehensive ablation study. Removing the reflection and scattering masks leads to a substantial performance drop on the Non-Lambertian dataset [4], confirming that decoupling non-Lambertian effects is important for accurate depth recovery. Furthermore, disabling the phase regularization term results in noisy depth maps around specular highlights, highlighting the necessity of enforcing smooth phase transitions in diffuse regions. These results collectively demonstrate that each module contributes significantly to the overall superiority of the Unified Dual-Mask Physical Model. With these experimental results establishing the framework's effectiveness, we now turn to summarizing the core contributions and broader implications of this work.